\documentclass[11pt,a4paper]{article}

% ── Codificación y fuentes ────────────────────────────────────────────────────
\usepackage[utf8]{inputenc}
\usepackage[T1]{fontenc}
\usepackage{lmodern}
\usepackage[spanish]{babel}

% ── Geometría y espaciado ─────────────────────────────────────────────────────
\usepackage[top=2.5cm, bottom=2.5cm, left=2.8cm, right=2.8cm]{geometry}
\usepackage{setspace}
\setstretch{1.15}
\usepackage{parskip}

% ── Tablas ────────────────────────────────────────────────────────────────────
\usepackage{booktabs}
\usepackage{tabularx}
\usepackage{array}
\usepackage{longtable}
\usepackage{enumitem}

% ── Colores y cajas ───────────────────────────────────────────────────────────
\usepackage[dvipsnames]{xcolor}
\usepackage{tcolorbox}
\tcbuselibrary{skins, breakable}

% ── Cabeceras y pies de página ────────────────────────────────────────────────
\usepackage{fancyhdr}
\usepackage{lastpage}
\pagestyle{fancy}
\fancyhf{}
\fancyhead[L]{\small\textbf{Informe de Evaluación} --- Física para Bioanalistas}
\fancyhead[R]{\small camila cabrera 32937064}
\fancyfoot[C]{\small Página \thepage\ de \pageref{LastPage}}
\renewcommand{\headrulewidth}{0.4pt}

% ── Hiperenlaces ──────────────────────────────────────────────────────────────
\usepackage[colorlinks=true, linkcolor=NavyBlue, urlcolor=NavyBlue]{hyperref}

% ── Paleta de colores personalizados ─────────────────────────────────────────
\definecolor{desempeno_excelente}{HTML}{1A6B2F}
\definecolor{desempeno_bueno}{HTML}{2E5A9C}
\definecolor{desempeno_suficiente}{HTML}{8B6914}
\definecolor{desempeno_deficiente}{HTML}{C0392B}
\definecolor{desempeno_insuficiente}{HTML}{7B241C}
\definecolor{grisFondo}{HTML}{F5F5F5}
\definecolor{azulTitulo}{HTML}{1B2A6B}
\definecolor{verdeFortaleza}{HTML}{1A6B2F}
\definecolor{naranjaMejora}{HTML}{A04000}

% ── Estilos de cajas ─────────────────────────────────────────────────────────
\tcbset{
  cajaTitulo/.style={
    enhanced, breakable,
    colback=azulTitulo!8, colframe=azulTitulo,
    fonttitle=\bfseries\large, coltitle=white,
    attach boxed title to top left={yshift=-2mm, xshift=4mm},
    boxed title style={colback=azulTitulo},
    top=6pt, bottom=6pt, left=8pt, right=8pt,
  },
  cajaFortaleza/.style={
    enhanced, breakable,
    colback=verdeFortaleza!6, colframe=verdeFortaleza!60,
    fonttitle=\bfseries, coltitle=verdeFortaleza,
    top=4pt, bottom=4pt, left=8pt, right=8pt,
  },
  cajaMejora/.style={
    enhanced, breakable,
    colback=naranjaMejora!6, colframe=naranjaMejora!60,
    fonttitle=\bfseries, coltitle=naranjaMejora,
    top=4pt, bottom=4pt, left=8pt, right=8pt,
  },
  cajaRecomendacion/.style={
    enhanced, breakable,
    colback=grisFondo, colframe=gray!50,
    fonttitle=\bfseries, coltitle=black,
    top=4pt, bottom=4pt, left=8pt, right=8pt,
  },
}

% ─────────────────────────────────────────────────────────────────────────────
\begin{document}

% ══ PORTADA ══════════════════════════════════════════════════════════════════
\begin{center}
  {\Large\bfseries\color{azulTitulo} Universidad de Oriente/Núcleo Sucre --- Física para Bioanalistas}\\[4pt]
  {\large Informe de Evaluación Preliminar}\\[12pt]
  \rule{\linewidth}{1.2pt}\\[6pt]
  {\LARGE\bfseries camila cabrera 32937064}\\[4pt]
  \rule{\linewidth}{0.6pt}
\end{center}

% ══ FICHA TÉCNICA ════════════════════════════════════════════════════════════
\vspace{4pt}
\begin{tcolorbox}[cajaTitulo, title=Datos del informe]
\begin{tabular}{@{}llll@{}}
  \textbf{Estudiante:}  & camila cabrera 32937064  &
  \textbf{Fecha:}       & 2026-06-25 \\[4pt]
  \textbf{Archivo PDF:} & \texttt{camila\_cabrera\_32937064.pdf}  &
  \textbf{Páginas:}     & 3 \\
\end{tabular}
\end{tcolorbox}

% ══ RESULTADO GLOBAL ═════════════════════════════════════════════════════════
\vspace{6pt}
\begin{center}
  \begin{tcolorbox}[
    enhanced,
    colback=white, colframe=desempeno_excelente,
    width=0.62\linewidth,
    boxrule=2pt,
    halign=center, valign=center,
    top=8pt, bottom=8pt,
  ]
    \centering
    {\large\bfseries Puntaje sugerido}\\[6pt]
    {\Huge\bfseries\color{desempeno_excelente} 9.8/10}\\[6pt]
    {\large Nivel de desempeño: \textbf{\color{desempeno_excelente} Excelente}}
  \end{tcolorbox}
\end{center}

% ══ RESUMEN GENERAL ══════════════════════════════════════════════════════════
\section*{Resumen general}
El estudiante demuestra un excelente dominio de los principios de la física de fluidos aplicados al bioanálisis. Ha resuelto correctamente la mayoría de los problemas, mostrando una sólida comprensión conceptual y habilidades procedimentales precisas. Su trabajo es claro, organizado y los cálculos son mayormente impecables.

% ══ FORTALEZAS Y ÁREAS DE MEJORA ════════════════════════════════════════════
\vspace{4pt}
\begin{minipage}[t]{0.48\linewidth}
  \begin{tcolorbox}[cajaFortaleza, title={Fortalezas}]
\begin{itemize}[leftmargin=1.5em, itemsep=2pt]
  \item Sólida comprensión de los principios de hidrostática, Arquímedes, continuidad, Bernoulli y Poiseuille.
  \item Habilidad para aplicar fórmulas correctamente y realizar cálculos precisos.
  \item Correcto manejo de unidades y conversiones (ej. mm a m).
  \item Organización clara y legible en el desarrollo de los problemas.
  \item Capacidad para interpretar correctamente las condiciones de los problemas (ej. variables constantes en análisis proporcional).
\end{itemize}
  \end{tcolorbox}
\end{minipage}
\hfill
\begin{minipage}[t]{0.48\linewidth}
  \begin{tcolorbox}[cajaMejora, title={Areas de mejora}]
\begin{itemize}[leftmargin=1.5em, itemsep=2pt]
  \item Precisión en la asignación inicial de variables para evitar confusiones, aunque no afectó el resultado final en el problema 1.
\end{itemize}
  \end{tcolorbox}
\end{minipage}

% ══ OBSERVACIONES POR PREGUNTA ═══════════════════════════════════════════════
\section*{Observaciones por pregunta}

\begin{longtable}{>{\bfseries}p{2.8cm} p{1.6cm} p{7.0cm} p{2.8cm}}
  \toprule
  \textbf{Pregunta} & \textbf{Puntaje} & \textbf{Evaluación} & \textbf{Errores detectados} \\
  \midrule
  \endhead
  \bottomrule
  \endfoot
    \textbf{Pregunta 1: Presión Hidrostática y Venosa} & 1.8/2.0 & El estudiante aplicó correctamente la fórmula de presión hidrostática (P = $\rho$gh) para calcular la altura. Realizó los cálculos de manera precisa y la conversión de unidades (metros a centímetros) fue correcta. Demuestra una buena comprensión del concepto. & \textit{Confusión inicial en la asignación de variables (escribió 'P = 1060 kg/m$^{3}$' antes de usar el valor correcto para la densidad en el cálculo, aunque luego usó la densidad correctamente).} \\
    \midrule
    \textbf{Pregunta 2: Principio de Arquímedes} & 2.0/2.0 & La resolución es impecable. El estudiante calculó correctamente la fuerza de flotación, el volumen del objeto y, finalmente, la densidad del tejido. Todos los pasos son lógicos, los cálculos precisos y las unidades consistentes. Excelente aplicación del principio de Arquímedes. & \textit{---} \\
    \midrule
    \textbf{Pregunta 3: Hemodinámica (Continuidad y Bernoulli)} & 2.0/2.0 & El estudiante aplicó correctamente tanto la ecuación de continuidad para determinar la velocidad en la zona estenosada, como la ecuación de Bernoulli para calcular la presión final. Los cálculos son exactos y el procedimiento es claro y bien estructurado. Demuestra un dominio completo de ambos principios. & \textit{---} \\
    \midrule
    \textbf{Pregunta 4: Ley de Poiseuille (Análisis Proporcional)} & 2.0/2.0 & La respuesta es correcta. El estudiante identificó que, bajo las condiciones dadas (diferencia de presión y viscosidad constantes), el caudal es directamente proporcional a la cuarta potencia del radio (Q $\propto$ r$^{4}$). Calculó correctamente la reducción porcentual del caudal al disminuir el radio en un 20\%. & \textit{---} \\
    \midrule
    \textbf{Pregunta 5: Flujo Viscoso y Resistencia} & 2.0/2.0 & El estudiante aplicó correctamente la Ley de Poiseuille para calcular la diferencia de presión requerida. Realizó las conversiones de unidades necesarias (mm a m) y los cálculos fueron precisos, incluyendo la potencia del radio (r$^{4}$). La respuesta final es correcta y bien expresada. & \textit{---} \\
    \midrule
\end{longtable}

% ══ ERRORES DETECTADOS ═══════════════════════════════════════════════════════
\section*{Análisis de errores}

\subsection*{Errores conceptuales}
\textit{Ninguno identificado.}

\subsection*{Errores procedimentales}
\textit{Ninguno identificado.}

% ══ RECOMENDACIONES AL ESTUDIANTE ════════════════════════════════════════════
\section*{Retroalimentación para el estudiante}
\begin{tcolorbox}[cajaRecomendacion, title={Dirigido a: camila cabrera 32937064}]
¡Excelente trabajo! Has demostrado una comprensión muy sólida de los conceptos de física de fluidos y una gran habilidad para aplicarlos en la resolución de problemas. Tu desarrollo es claro y tus cálculos son precisos. Solo un pequeño detalle a cuidar es la asignación inicial de variables para asegurar que siempre representen la magnitud correcta (por ejemplo, no confundir densidad con presión al inicio de un problema), aunque en este examen lo corregiste en el cálculo. Sigue así, tu desempeño es sobresaliente.
\end{tcolorbox}

% ══ NOTA PARA EL PROFESOR ════════════════════════════════════════════════════
\section*{Nota para el profesor}
\begin{tcolorbox}[
  enhanced, breakable,
  colback=yellow!8, colframe=yellow!60!black,
  fonttitle=\bfseries, coltitle=black,
  title={Observaciones para revisión manual},
  top=4pt, bottom=4pt, left=8pt, right=8pt,
]
El estudiante ha resuelto el examen con un desempeño excelente. Todos los problemas fueron abordados correctamente, mostrando una comprensión profunda de los principios físicos y una ejecución impecable. El único punto menor a señalar es una pequeña confusión inicial en la asignación de una variable en la Pregunta 1 (escribió 'P = 1060 kg/m$^{3}$' antes de usar 'P = 2500 Pa'), pero el cálculo posterior usó los valores correctos para densidad y presión, por lo que no afectó el resultado final. Se sugiere un puntaje de 9.8/10.
\end{tcolorbox}

% ══ PIE DE INFORME ════════════════════════════════════════════════════════════
\vspace{12pt}
\begin{center}
  \footnotesize\color{gray}
  Informe generado automáticamente por el sistema de evaluación con IA (gemini-2.5-flash) y soporte LaTex. \\
  Este documento es una evaluación \textbf{preliminar} para ser revisado por el profesor antes de ser entregado al 
  estudiante. \\
  Generado el 2026-06-25 \textbullet\ BIOANALISIS Sistema Académico de Evaluación.
  \end{center}

\end{document}
